\documentclass[11pt]{article}
\usepackage[margin=1.1in]{geometry}
\usepackage{amsmath,amssymb,amsthm}
\usepackage{booktabs}
\usepackage[colorlinks=true,linkcolor=blue,citecolor=blue,urlcolor=blue]{hyperref}

\newtheorem{theorem}{Theorem}[section]
\newtheorem{lemma}[theorem]{Lemma}
\newtheorem{proposition}[theorem]{Proposition}
\newtheorem{corollary}[theorem]{Corollary}
\newtheorem{conjecture}[theorem]{Conjecture}
\theoremstyle{definition}
\newtheorem{definition}[theorem]{Definition}
\newtheorem{remark}[theorem]{Remark}
\newtheorem{question}[theorem]{Question}

\newcommand{\eps}{\varepsilon}
\DeclareMathOperator{\lc}{LC}

\title{Unimodality of independence polynomials of hub spiders:\\
log-concavity certificates and a reduction to a single inequality}
\author{Jongyun (Max) Han\thanks{madmax0404@maxandomnis.com}}
\date{June 2026 --- draft v0.4}

\begin{document}
\maketitle

\begin{abstract}
A conjecture of Alavi, Malde, Schwenk, and Erd\H{o}s (1987) asserts that
the independent-set sequence of every tree and forest is unimodal; it is
problem~993 in the Erd\H{o}s problems collection.  We study the family of
\emph{hub spiders} $S(c_1,\dots,c_h)$: a root joined to $h$ hubs, the
$i$-th hub bearing $c_i$ pendant paths of length two.  This family
contains the trees $T_{a,b,c}$ of Kadrawi--Levit, the first known trees
whose independence sequence fails log-concavity, and for which
unimodality was previously known only in special cases.

We develop a reduction of unimodality for this family to a single
explicit inequality: the log-concavity in $t$ of the coefficient flow
$G(t)=[x^k]\prod_i\bigl((1+2x)^{c_i}+t\,x(1+x)^{c_i}\bigr)$ at band
positions $k$.  The route runs through: a quantitative
ultra-log-concavity theorem for the hub factor
$H_c=(1+2x)^c+x(1+x)^c$ with virtual degree $c+1+\lceil c/30\rceil$;
an exact identity expressing the log-concavity defect of $I(S)$
through the relative-mass sequence $\eps_k=(F_k-g_k)/g_k$; sharp
per-class growth bounds obtained from exact window-mean identities
and Jensen's inequality; and the concavity of the flow, which
controls all moments of the class-mass profile.  The flow concavity
is proved completely for $h=2$, for arbitrary arms --- by an
elementary argument (a trinomial collapse, symmetrization, AM--GM,
and Vandermonde) extended through a $t$-cancellation identity and a
coefficientwise Schur lemma; at
uniform multisets it is equivalent to a symmetric pairwise
inequality (the \emph{E-core}), which is proved in the limit
$c\to\infty$ for all $h$ and $k$ by a five-line calculus bound and
established for all integer $c$ on a large finite range of $(k,h)$
by an exact polynomial certification method.  Consequently the
unimodality of $I(S)$ is proved unconditionally for $h=1$ (all
$c$) and for $82$ further families: $55$ uniform ones including
$S(4^5)$, a particularly fragile uniform hub spider found by the search, and
the diagonal $T_{c,c,c}$ for $c\le 10$; eleven mixed-arm
two-hub members of the certificate grid; and sixteen three-hub
adjacent-two-value families including $T_{3,4,4}$ and $T_{3,3,4}$,
the first published non-log-concave trees.  Two natural pairwise strengthenings of
the flow concavity are shown to be false, with explicit hub-spider
witnesses; the reduction is stated in the exact form the
counterexamples leave standing.

We also record structural results from a systematic counterexample
search: a balance invariant that measures distance from unimodality
failure, an exact offset law for the location of log-concavity failures
in uniform hub spiders, and the observation that the
Levit--Mandrescu/Basit--Galvin decreasing zone protects this junction
mechanism from becoming a unimodality failure.
\end{abstract}

\section{Introduction}\label{sec:intro}

For a finite graph $G$, write $i_k(G)$ for the number of independent
sets of size $k$, $\alpha=\alpha(G)$ for the independence number, and
\[
  I(G;x)\;=\;\sum_{k=0}^{\alpha} i_k(G)\,x^k
\]
for the independence polynomial.  A sequence $(a_k)$ is \emph{unimodal}
if $a_0\le\cdots\le a_m\ge\cdots\ge a_\alpha$ for some $m$, and
\emph{log-concave} (LC) if $a_k^2\ge a_{k-1}a_{k+1}$ for all internal
$k$; for positive sequences log-concavity implies unimodality.
Alavi, Malde, Schwenk, and Erd\H{o}s~\cite{AMSE87} proved that the
independence sequence of a general graph can be arbitrarily far from
unimodal and asked whether it is unimodal for every tree and forest;
the question appears as problem~993 of~\cite{ErdosProblems}.

The conjecture has stimulated a substantial literature.  Levit and
Mandrescu proved that for K\H{o}nig--Egerv\'ary graphs (hence all
forests) the final third of the sequence is nonincreasing; Basit and
Galvin~\cite{BasitGalvin} extended the decreasing zone and proved that
random labelled trees are almost surely ``almost unimodal''.
Galvin and Hilyard~\cite{GalvinHilyard} established unimodality for
several recursively structured tree families.  On the negative side,
Kadrawi and Levit~\cite{KadrawiLevit} gave infinite families extending
the $26$-vertex examples whose independence sequences are not
log-concave.  Galvin~\cite{GalvinNonLC} later constructed families with
log-concavity failures at unbounded independence number; machine
search~\cite{RamosSun} produced many further examples; Li~\cite{Li}
proved unimodality for two infinite Kadrawi--Levit families; and
Bautista-Ramos, Guill\'en-Galv\'an and G\'omez-Salgado~\cite{BautistaRamos}
developed linear recurrences and further non-log-concave families.  The
hub-spider subfamily below isolates the junction mechanism in the
Kadrawi--Levit examples and in many examples found by the search.

\begin{definition}[Hub spiders]\label{def:hub}
For $h\ge 1$ and $c_1,\dots,c_h\ge 1$, the \emph{hub spider}
$S(c_1,\dots,c_h)$ is the tree consisting of a root $r$, hubs
$v_1,\dots,v_h$ adjacent to $r$, and, for each $i$, $c_i$ pendant paths
of length two attached to $v_i$.  It has $n=1+h+2C$ vertices and
independence number $D=h+C$, where $C=\sum_i c_i$.  We write $S(c^h)$
in the uniform case $c_1=\cdots=c_h=c$.
\end{definition}

The family contains the Kadrawi--Levit trees $T_{a,b,c}=S(a,b,c)$,
including the unstarred families treated by Li~\cite{Li}, and captures
many hub-spider examples produced by the machine search~\cite{RamosSun}.
By the root
decomposition,
\begin{equation}\label{eq:PFG}
  I(S(c_1,\dots,c_h);x) \;=\; P \;=\; F + x\,g,\qquad
  F=\prod_{i=1}^h H_{c_i},\quad g=(1+2x)^{C},
\end{equation}
where $H_c=(1+2x)^c+x(1+x)^c$ is the independence polynomial of a hub
with $c$ pendant $P_2$'s, the two summands according to whether the hub
is excluded or included.

\subsection{Results}

Our first group of results is unconditional.

\begin{theorem}[Unconditional scope]\label{thm:scope}
$I(S(c_1,\dots,c_h);x)$ has a unimodal coefficient sequence in the
following cases:
\begin{enumerate}
  \item $h=1$, all $c\ge 1$ (indeed $I$ is log-concave);
  \item the uniform families $S(c^h)$ with
    \[
      \begin{array}{c|cccccccccc}
        h & 1 & 2 & 3 & 4 & 5 & 6 & 7 & 8 & 9{-}10\\\hline
        c\le & \text{all} & 10 & 10 & 10 & 7 & 6 & 5 & 3 & 2
      \end{array}
    \]
  \item the eleven mixed-arm spiders of the certificate grid:
    $S(1,8)$, $S(2,4)$, $S(3,4)$, $S(4,7)$, $S(4,9)$, $S(4,11)$,
    $S(6,8)$, $S(6,12)$, $S(7,8)$, $S(10,12)$, $S(11,12)$;
  \item the sixteen three-hub adjacent-two-value spiders
    $S(a^p,(a{+}1)^{3-p})$, $1\le a\le8$, $p\in\{1,2\}$ ---
    including $T_{3,4,4}=S(3,4,4)$ and $T_{3,3,4}=S(3,3,4)$, the
    first published non-log-concave trees.
\end{enumerate}
In particular $S(4^5)$, a particularly fragile uniform hub spider found
by the search (Section~\ref{sec:extremal}), and the uniform Kadrawi--Levit
diagonal $T_{c,c,c}$, $c\le 10$, have unimodal independence sequences.
\end{theorem}

The proof of Theorem~\ref{thm:scope} is a composition of human-proved
lemmas with machine-verified components, all executed in exact
rational arithmetic and reproducible from the public scripts; the
methodology is described in Section~\ref{sec:verification}.

The structural heart of the paper is a reduction theorem.

\begin{theorem}[Reduction to the flow concavity]\label{thm:reduction}
Fix $h\ge 2$ and $c_1,\dots,c_h$, and let $b_i=x(1+x)^{c_i}$,
$E_i=(1+2x)^{c_i}$.  Suppose that
\begin{enumerate}
\item[(i)] the certificate inequalities of
  Section~\ref{sec:verification} hold for $(c_1,\dots,c_h)$ at every
  band position (verified exactly for every multiset in the grid
  $\mathcal G$ of Section~\ref{sec:verification}, and conjectured in
  general, Conjecture~\ref{conj:g2}); and
\item[(ii)] at every position $k$ in the band \eqref{eq:band}, the
  demotion flow $G(t)=\bigl[\prod_{l=1}^h(E_l+t\,b_l)\bigr]_k$ is
  log-concave in $t$ on $[0,2]$:
\begin{equation}\label{eq:flowcc}
  G(t)\,G''(t)\;\le\;G'(t)^2 ,\qquad 0\le t\le 2 .
\end{equation}
\end{enumerate}
Then $I(S(c_1,\dots,c_h);x)$ is unimodal.
\end{theorem}

Expanding \eqref{eq:flowcc} by the product rule
(Proposition~\ref{prop:dd}), hypothesis (ii) is the aggregate
inequality $\sum_{i\ne j}(G\,y_{ij}-x_ix_j)\le\sum_ix_i^2$, where
$A_l=E_l+t\,b_l$, $x_i=[b_i\prod_{l\ne i}A_l]_k$ and
$y_{ij}=[b_ib_j\prod_{l\ne i,j}A_l]_k$.  At uniform multisets the
diagonal flow makes all $x_i$ equal and all $y_{ij}$ equal, and (ii)
becomes the \emph{symmetric E-core}
\begin{equation}\label{eq:sym}
  [b_ib_jW]_k\,[A_iA_jW]_k\;\le\;
  \Bigl(1+\frac1{h-1}\Bigr)[b_iA_jW]_k\,[b_jA_iW]_k,
  \qquad W=\prod_{l\ne i,j}A_l,
\end{equation}
which is the form all certifications below address.

We caution --- twice --- against strengthening (ii) to a
\emph{pairwise} statement for general multisets; both natural
pairwise forms were intermediate conjectures of this project,
refuted by adversarial scanning.  The symmetric form \eqref{eq:sym}
fails at lopsided pairs: at $(1,c,1^{h-2})$ with all-$E$ background
it fails in band from $c\approx16$, the normalized defect reaching
$4.35$ at $(h,c)=(8,64)$.  The weighted form
$G\,y_{ij}-x_ix_j\le(x_i^2+x_j^2)/(2(h-1))$ --- the exact-allocation
split, which restricts to \eqref{eq:sym} at $x_i=x_j$ --- fails at
$1$-heavy backgrounds: its smallest violation is the $107$-vertex
spider $S(1,27,1^{16})$ at $k=5$, and at $(128,128,1^{192})$, $k=7$
it fails by the factor $2.53$.  The failure mode is an allocation
artifact: the uniform $1/(h-1)$ budget starves the heavy pair while
the many light pairs underdraw.  The aggregate (ii) itself passes
every exact scan --- $14{,}410$ evaluations across the flow
parameter, including all $200$ pairwise-violation witnesses, with
worst ratio $0.085$ and margins that \emph{grow} (to $-5.3$ at
background size $1024$) exactly where the pairwise splits die
(Section~\ref{sec:verification}).

Four results establish \eqref{eq:flowcc} unconditionally on large
families.

\begin{theorem}[E-core at $h=2$]\label{thm:h2}
For $h=2$ the symmetric form holds with constant $2$ for all $c\ge1$
and all $k\ge0$, with no band restriction; by AM--GM this implies
the flow concavity \eqref{eq:flowcc} at $h=2$, $c_1=c_2$.
\end{theorem}

The proof (Section~\ref{sec:p6}) is elementary and short: the bracket
$[bE]_k$ collapses to a coefficient of the trinomial power
$(1+3x+2x^2)^c$, and a symmetrization, the AM--GM inequality, the
Vandermonde identity, and log-concavity of binomial coefficients finish
in four lines.

\begin{theorem}[Flow concavity at $h=2$, arbitrary arms]\label{thm:h2mixed}
For $h=2$ with any arms $(c_1,c_2)$, the flow concavity
\eqref{eq:flowcc} holds at every position $k\ge0$, with no band
restriction.  In particular hypothesis (ii) of
Theorem~\ref{thm:reduction} holds for every two-hub spider.
\end{theorem}

The proof (Section~\ref{sec:p6}) combines two new lemmas with the
core of Theorem~\ref{thm:h2}: a \emph{$t$-cancellation} lemma --- for
any fixed background the pairwise defect $G y_{ij}-x_ix_j$ is
independent of the flow parameter on the two pair slots, a unipotent
congruence of the $2\times2$ bracket matrix --- and a
coefficientwise Schur lemma: among the splits $c_1+c_2=s$, the
coefficients of $(1+x)^{c_1}(1+2x)^{c_2}+(1+x)^{c_2}(1+2x)^{c_1}$ are
minimized at the balanced one.  For even $s$ the balanced split is
the core of Theorem~\ref{thm:h2}; odd $s$ needs one further
inequality (the \emph{odd core}, Lemma~\ref{lem:oddcore}), whose
four-line proof runs on Pascal's rule, log-concavity of a binomial
row, and the numerical accident $\tfrac32>\sqrt2$.

\begin{theorem}[Limit calculus]\label{thm:limit}
For each fixed $k$ and $h\ge3$, the defect of the symmetric form in
the uniform case with all-$E$ slots and background (the binding
configuration in all sweeps) converges as $c\to\infty$ to
\[
  \Delta_\infty(k,\ell)=\Bigl(1-\tfrac1k\Bigr)
  \Bigl(1+\tfrac1{2\ell+2}\Bigr)^{2}
  \Bigl(1-\tfrac1{(2\ell+3)^2}\Bigr)^{k}-1,
  \qquad \ell=h-2,
\]
and $(h-1)\sup_{k\ge2}\Delta_\infty(k,h-2)\le 1.1/(2h-1)\le 0.22$
for every $h\ge 3$: the limit is uniformly at least $4.5$ times below
the required bound.
\end{theorem}

\begin{theorem}[Finite certification]\label{thm:cert}
The symmetric form \eqref{eq:sym} (hence, at uniform multisets, the
flow concavity \eqref{eq:flowcc}) holds, for all integer
$c\ge1$ in the uniform case, at every $(k,h)$ with $3\le h\le 7$, $k\le 28$ and
$8\le h\le 10$, $k\le 20$, and every background pattern: $996$
instances, each certified by an exact polynomial argument
(Section~\ref{sec:verification}).
\end{theorem}

Finally, the counterexample search that initiated this work produced
extremal information of independent interest
(Section~\ref{sec:extremal}): an exact \emph{offset law} --- in
$S(c^h)$ with $c\ge4$ the unique log-concavity failure sits exactly
$h-2$ positions below the top degree --- and the \emph{record
fragility} tree, a $48$-vertex junction tree (four hub gadgets and
one exotic gadget) whose log-concavity
\emph{balance} (Definition~\ref{def:balance}) is $0.0451$, more than
twice the previous literature record, yet still $22$ times short of
what a unimodality counterexample requires.

\subsection{Method overview}\label{sec:overview}

Write $P=F+xg$ as in \eqref{eq:PFG} and let
\[
  \begin{gathered}
  \eps_k=\frac{F_k-g_k}{g_k},\qquad d_k=\eps_{k+1}-\eps_k,\\
  \rho_k=\frac{g_k}{g_{k-1}}=\frac{2(C-k+1)}{k},\qquad
  \lambda_k=\frac{(k-1)(C-k)}{(k+1)(C-k+2)} .
  \end{gathered}
\]
The engine of the reduction is the exact identity
(Lemma~\ref{lem:epsid})
\[
  \underbrace{2F_kg_{k-1}-F_{k-1}g_k-F_{k+1}g_{k-2}}_{\text{defect}(k)}
  \;=\; g_kg_{k-1}\Bigl[(1-\lambda_k)(1+\eps_k)+d_{k-1}-\lambda_kd_k\Bigr],
\]
which converts band log-concavity of $P$ into a one-dimensional
statement about the sequence $(\eps_k)$, absorbed against the
Newton-strength margin $\lc_F(k)\ge F_k^2\,(N_F+1)/\bigl((k+1)(N_F-k+1)\bigr)$
provided by the virtual ultra-log-concavity of the factors
(Theorem~\ref{thm:a2}).  Class decomposition of $F-g$, exact
window-mean identities, and two Jensen arguments give closed-form
two-sided bounds on every class growth; the \emph{demotion flow}
$\phi(t)=\ln[x^k]\prod_i(E_i+tb_i)$ supplies all moment bounds for the
class-mass profile via its concavity --- which is exactly the single
inequality \eqref{eq:flowcc} of Theorem~\ref{thm:reduction}, and which
at uniform multisets is equivalent to the symmetric E-core
\eqref{eq:sym}.  The final assembly into explicit certificates,
verified in exact rational arithmetic, is described in
Section~\ref{sec:verification}.

\section{Setup and notation}\label{sec:setup}

Throughout, $[u]_k$ (or $u_k$, or $u[k]$) denotes the coefficient of
$x^k$ in $u$; $\lc_u(k)=u_k^2-u_{k-1}u_{k+1}$; $r^u(j)=u_j/u_{j-1}$.
Fixed symbols: $C=\sum_ic_i$, $D=h+C$, $n=1+h+2C$,
$g=(1+2x)^C$, $\rho_k=2(C-k+1)/k$,
$\lambda_k=(k-1)(C-k)/\bigl((k+1)(C-k+2)\bigr)$,
$\eps_k=(F_k-g_k)/g_k$, $d_k=\eps_{k+1}-\eps_k$,
$N_c=c+1+\lceil c/30\rceil$, $N_F=\sum_iN_{c_i}$,
$m_k=(N_F+1)/\bigl((k+1)(N_F-k+1)\bigr)$,
$q_k=k/(2(C-k+1))$, $\sigma_k=\rho_k/\rho_{k+1}-1$.

For \eqref{eq:PFG} write $b_i=x(1+x)^{c_i}$, so $H_{c_i}=E_i+b_i$.
Expanding the product over the choice of summand in each factor,
\[
  F-g \;=\; R\;=\;\sum_{\emptyset\ne S\subseteq[h]} t_S,\qquad
  t_S=x^{|S|}\prod_{i\in S}(1+x)^{c_i}\prod_{i\notin S}(1+2x)^{c_i},
\]
and we write $W_s(k)=\bigl(\sum_{|S|=s}t_S[k]\bigr)/g_k$ for the class
masses, $\eps_k=\sum_s W_s(k)$.

The \emph{band} is
\begin{equation}\label{eq:band}
  2\;\le\;k\;\le\;k_A:=\min(\ell_{\mathrm{BG}}-1,\;k_0-1,\;C-1),
\end{equation}
where $k_0=\lceil(2D-1)/3\rceil$ (Levit--Mandrescu) and
$\ell_{\mathrm{BG}}=\lceil D(n-1)/(D+n)\rceil$ (Basit--Galvin), and set
$k_{\mathrm{dec}}=\min(k_0,\ell_{\mathrm{BG}})$.  Beyond
$k_{\mathrm{dec}}$ the sequence is nonincreasing by those theorems, and
positions between $C$ and the decreasing zone are handled separately
(Section~\ref{sec:assembly}).

\section{The composition skeleton}\label{sec:skeleton}

\begin{lemma}[Seam]\label{lem:seam}
Let $P$ be a positive sequence of degree $D$, log-concave at every
$k\in[1,k_{\mathrm{dec}}-1]$, and nonincreasing on
$[k_{\mathrm{dec}},D]$.  Then $P$ is unimodal.
\end{lemma}

\begin{proof}
Log-concavity at $k$ means $r_k\ge r_{k+1}$ for $r_k=P_k/P_{k-1}$, so
$r_1\ge\cdots\ge r_{k_{\mathrm{dec}}}$ and the ratio sequence crosses
$1$ at most once on that range; hence $P$ is unimodal on
$[0,k_{\mathrm{dec}}]$, with mode $m$ say.  On $[m,k_{\mathrm{dec}}]$
the ratios are $\le1$, so $P$ is nonincreasing there, and it continues
nonincreasing on $[k_{\mathrm{dec}},D]$ by hypothesis.
\end{proof}

The case $k=1$ of the log-concavity hypothesis is universal: for any
tree on $n$ vertices, $i_1=n$ and $i_2=\binom n2-(n-1)$, and
$n^2\ge\binom n2-(n-1)$ reads $n^2+3n-2\ge0$.

\subsection{Band B}\label{sec:bandB}

When the multiset is leaf-heavy (roughly $h\gtrsim 0.62\,C$) the
decreasing zone begins above $C$, leaving positions
$k\in[C,k_{\mathrm{dec}}-1]$ outside the band \eqref{eq:band}.  These
are handled directly.  Since $xg$ has support $[1,C+1]$, we have
$P_k=F_k$ for $k\ge C+2$; hence for $k\ge C+3$ all three coefficients
in $\lc_P(k)$ are pure $F$ and $\lc_P(k)=\lc_F(k)\ge0$ by
Lemma~\ref{lem:a0} and Lemma~\ref{lem:a1}.  At the three remaining
positions $k\in\{C,C+1,C+2\}$ the corrections to $\lc_F(k)$ are
explicit in the top three coefficients of $g$:
\[
  \lc_P(C{+}2)-\lc_F(C{+}2)=-g_CF_{C+3},\qquad
  \lc_P(C{+}1)-\lc_F(C{+}1)=2F_{C+1}g_C+g_C^2-g_{C-1}F_{C+2},
\]
and similarly at $k=C$ with $g_{C-2},g_{C-1},g_C$.  The negative
terms are absorbed by the margin of Theorem~\ref{thm:a2} using the
explicit mass bound
\begin{equation}\label{eq:W1C}
  \eps_C\;\ge\;W_1(C)\;=\;\sum_{i=1}^h\frac{C+c_i}{2^{\,c_i+1}},
\end{equation}
which follows from the top two coefficients of
$u_i=(1+x)^{c_i}(1+2x)^{C-c_i}$, namely $u_i[C]=2^{C-c_i}$ and
$u_i[C-1]=2^{C-c_i-1}(C+c_i)$.  In the leaf-heavy regime
\eqref{eq:W1C} is large (each $c_i\le2$ term contributes at least
$(C+1)/8$), so the $g$-tail is negligible against $F$.

The explicit threshold is $C\ge55$.  Put $A_c=2^c+c$, the
one-below-top coefficient of $H_c$.  Since Band~B is nonempty only
when $h^2+Ch-C^2-C>0$, we have $h\ge3C/5$, and if $r$ is the number
of indices with $c_i=1$, then $r\ge2h-C$.  The reversed top expansion
of $F$ gives
\[
  F_C\ge e_h(A_{c_i}),\qquad
  F_{C+1}\ge e_{h-1}(A_{c_i}),\qquad
  F_{C+2}\ge e_{h-2}(A_{c_i}).
\]
The compression inequality
\[
  (2^a+a)(2^b+b)\ge 3(2^{a+b-1}+a+b-1),\qquad a,b\ge1,
\]
therefore gives
$F_C\ge 3^{h-1}(2^{C-h+1}+C-h+1)$; and $A_1=3$, $A_c\ge2^c$ give
\[
  F_{C+1}\ge r3^{r-1}2^{C-r},\qquad
  F_{C+2}\ge \binom r2 3^{r-2}2^{C-r}.
\]
Together with $N_F\le2C+h$ and
$r^F(k{+}1)\le(2C+h)(N_F-k)/((k{+}1)N_F)$ from
Theorem~\ref{thm:a2}, these bounds make the three absorptions at
$k=C,C+1,C+2$ explicit rational-exponential inequalities.  For
$C\ge55$ a quotient check shows that the five residuals are minimized
on the Band~B range at $(C,h)=(55,35)$, where all are positive.  The
remaining finite list is exactly the leaf-heavy multisets with
$C\le54$; these
are the $24{,}394$ integer partitions of $C-h$ over all admissible
$(C,h)$, comprising $54{,}049$ Band~B coefficient positions.  The
verifier checks each position in exact integer arithmetic and finds
zero failures.

\section{Log-concavity of the factors}\label{sec:factors}

\begin{lemma}[A0]\label{lem:a0}
For every $c\ge1$ and $t\in[0,2]$ the polynomial $E_c+t\,b_c$ has a
log-concave coefficient sequence.  Quantitatively, for $H_c=E_c+b_c$,
\[
  H_k^2-H_{k-1}H_{k+1}\;\ge\;
  \frac{4^k\binom ck^2(c+1)}{(k+1)(c-k+1)}
  \;-\;2^{k-1}\binom ck\binom c{k-1}\;\ge\;0 .
\]
\end{lemma}

\begin{proof}
With $A_k=\binom ck2^k$ and $B_k=\binom c{k-1}$ decompose
$\Delta_k=(A+tB)_k^2-(A+tB)_{k-1}(A+tB)_{k+1}$ into the two pure
margins and a cross term:
\[
  \Delta_k=\underbrace{\frac{4^k\binom ck^2(c+1)}{(k+1)(c-k+1)}}_{A\text{-margin}}
  +\;t^2\underbrace{\frac{\binom c{k-1}^2(c+1)}{k(c-k+2)}}_{B\text{-margin}}
  +\;t\,2^{k-1}\Bigl[3\tbinom ck\tbinom c{k-1}-4\tbinom c{k+1}\tbinom c{k-2}\Bigr].
\]
Both margins are exact evaluations of binomial log-concavity.  Since
$\binom c{k+1}\binom c{k-2}<\binom ck\binom c{k-1}$ always, the cross
deficit is at most $t\,2^{k-1}\binom ck\binom c{k-1}$, and for $t\le2$
the $A$-margin absorbs it because
$2^{k}(c+1)\ge k(k+1)$, which holds as $2^k\ge k$ and $c+1\ge k+1$.
\end{proof}

\begin{theorem}[A2: virtual ultra-log-concavity]\label{thm:a2}
$H_c$ is $N_c$-ULC with $N_c=c+1+\lceil c/30\rceil$: the sequence
$H_k/\binom{N_c}k$ is log-concave.  Consequently, by Liggett's
convolution theorem~\cite{Liggett}, $F=\prod_iH_{c_i}$ is $N_F$-ULC
with $N_F=\sum_iN_{c_i}$, and for all $k$
\[
  \lc_F(k)\;=\;F_k^2-F_{k-1}F_{k+1}\;\ge\;
  F_k^2\,\frac{N_F+1}{(k+1)(N_F-k+1)} .
\]
\end{theorem}

\begin{remark}
$H_c$ is not real-rooted for $c\ge3$ and not plainly ULC for
$c\ge41$, so a virtual degree is necessary; $N_c$ is nearly optimal
(empirically the minimal $N$ is $c+1$ for $c\le40$ and
${\sim}1.032c$ asymptotically).
\end{remark}

\begin{proof}
Write $C_j=\binom cj$, $A_k=C_k2^k$, $B_k=C_{k-1}$, $N=N_c$, and
$\Delta=H_k^2\,k(N-k)-H_{k-1}H_{k+1}(k+1)(N-k+1)$; we must show
$\Delta\ge0$ for $1\le k\le c$.  Three exact evaluations, using
$C_{k-1}C_{k+1}/C_k^2=k(c-k)/\bigl((k+1)(c-k+1)\bigr)$ and the
telescoping $(N-k)(c-k+1)-(N-k+1)(c-k)=N-c$:
\begin{align*}
  A\text{-part:}\quad
  &A_k^2k(N-k)-A_{k-1}A_{k+1}(k+1)(N-k+1)
   =2^kC_kC_{k-1}\cdot2^k(N-c);\\
  B\text{-part:}\quad
  &B_k^2k(N-k)-B_{k-1}B_{k+1}(k+1)(N-k+1)
   =\frac{C_{k-1}^2}{k(c-k+2)}
   \Bigl[k^2(N-c-1)\\
  &\hspace{8em} +(c-k+1)(N-k+1)\Bigr]\ge0;\\
  \text{cross:}\quad
  &2A_kB_kk(N-k)-(A_{k-1}B_{k+1}+A_{k+1}B_{k-1})(k+1)(N-k+1)\\
  &\qquad=2^kC_kC_{k-1}
  \Bigl[2k(N-k)-(k+1)(N-k+1)\bigl(\tfrac12+2\theta_k\bigr)\Bigr],\\
  &\qquad\theta_k=\tfrac{(k-1)(c-k)}{(k+1)(c-k+2)}.
\end{align*}
Dropping the nonnegative $B$-part and clearing denominators, it
suffices that
\[
  S(k,c,N):=2(c-k+2)\bigl[2^k(N-c)+2k(N-k)\bigr]
  -(N-k+1)\bigl[(k+1)(c-k+2)+4(k-1)(c-k)\bigr]\;\ge\;0 .
\]
$S$ is nondecreasing in $N$ (the $N$-derivative is
$2(c-k+2)(2^k+2k)-[(k+1)(c-k+2)+4(k-1)(c-k)]\ge0$ since
$2^{k+1}\ge k+1$ and $4k(c-k+2)\ge4(k-1)(c-k)$), so it is enough to
prove $S\ge0$ at the real value $N=c+1+c/30$.  For $k\le12$,
$P_k(c):=30\,S(k,c,c+1+c/30)$ is a quadratic in $c$ with positive
leading coefficient and $P_k(k)\ge0$, $P_k'(k)\ge0$ --- for example
$P_1=6(c+1)(11c+10)$, $P_2=c(39c+248)$, \dots,
$P_{12}=7913c^2+170618c-2496600$ --- hence $P_k\ge0$ on $c\ge k$.
For $k\ge13$, use $(k+1)(c-k+2)+4(k-1)(c-k)\le5(k+1)(c-k+2)$,
$N-k+1\le N+1\le\tfrac{31}{30}(c+30)$, and
$N-c\ge\tfrac{1}{30}(c+30)$:
\[
  S\;\ge\;\frac{(c-k+2)(c+30)}{30}\bigl[2^{k+1}-155(k+1)\bigr]\;\ge\;0,
\]
since $2^{k+1}\ge155(k+1)$ for $k\ge13$.  (The inequality $S\ge0$ was
additionally verified exactly for every $c\le3000$ and every $k\le
c$.)  The Liggett corollary follows since zero-padding a degree-$d$
$N$-ULC sequence to virtual degree $N$ preserves log-concavity of the
normalized sequence.
\end{proof}

\begin{lemma}[Closure]\label{lem:a1}
The convolution of positive log-concave sequences with no internal
zeros is log-concave.  Quantitatively,
$\lc_{u*v}(k)\ge\sum_j\lc_u(k-1-j)\,\lc_v(j+1)$.
\end{lemma}

\begin{proof}[Proof sketch]
The qualitative statement is classical (positive LC sequences with no
internal zeros are P\'olya frequency of order 2, and PF$_2$ is closed
under convolution; see e.g.~\cite{Karlin}).  The quantitative form
follows from the Cauchy--Binet expansion of the $2\times2$ minors of
the convolution into sums of products of paired spread minors of $u$
and $v$, of which the displayed sum keeps the adjacent pairs; it is
stated for completeness but only the qualitative statement is used in
this paper (the quantitative margins come from Theorem~\ref{thm:a2}).
\end{proof}

\begin{lemma}[D1: the case $h=1$]\label{lem:d1}
For every $c\ge1$, $I(S(c))=H_c+xE_c$ is log-concave; in fact the
mixed defect $2H_kg_{k-1}-H_{k-1}g_k-H_{k+1}g_{k-2}$ is nonnegative
for all $1\le k\le c$, where $g=E_c$.
\end{lemma}

\begin{proof}
Split $H=g+b$, $b_k=\binom c{k-1}$; the defect is linear in its first
argument, $\mathrm{defect}=W(g,g)+W(b,g)$.  The pure part has the
exact closed form
$W(g,g)(k)=g_kg_{k-1}-g_{k+1}g_{k-2}
=g_kg_{k-1}(2c+2)/\bigl((k+1)(c-k+2)\bigr)>0$, and the mixed part
evaluates to
$W(b,g)(k)=2^{k-2}\bigl[4\binom c{k-1}^2-5\binom c{k-2}\binom ck\bigr]
\ge-2^{k-2}\binom c{k-1}^2$ by log-concavity of the binomial row.
Absorption: $g_kg_{k-1}=\binom ck\binom c{k-1}2^{2k-1}$, so
$W(g,g)\ge2^{k-2}\binom c{k-1}^2$ reduces to
$2^{k+2}(c+1)(c-k+1)\ge k(k+1)(c-k+2)$, which holds since
$c-k+1\ge(c-k+2)/2$ and $2^{k+1}(c+1)\ge k(k+1)$.  Log-concavity of
$I(S(c))$ then follows from Lemma~\ref{lem:a0}, log-concavity of the
shifted binomial $xg$, and $\lc_P=\lc_H+\lc_{xg}+\mathrm{defect}$.
\end{proof}

\section{The \texorpdfstring{$\eps$}{eps}-calculus and the master
inequality}\label{sec:eps}

\begin{lemma}[The $\eps$-identity]\label{lem:epsid}
With the notation of Section~\ref{sec:overview},
\[
  2F_kg_{k-1}-F_{k-1}g_k-F_{k+1}g_{k-2}
  = g_kg_{k-1}\bigl[(1-\lambda_k)(1+\eps_k)+d_{k-1}-\lambda_kd_k\bigr].
\]
Moreover $\lc_P(k)\ge \mathrm{defect}(k)+\lc_F(k)$, and with the margin
of Theorem~\ref{thm:a2} the band log-concavity of $P$ follows from the
\emph{master inequality}
\begin{equation}\label{eq:master}
  \lambda_kd_k-d_{k-1}\;\le\;(1-\lambda_k)(1+\eps_k)
  +\rho_k m_k(1+\eps_k)^2,
  \qquad m_k=\frac{N_F+1}{(k+1)(N_F-k+1)} .
\end{equation}
\end{lemma}

\begin{proof}
Substituting $F_j=(1+\eps_j)g_j$ and dividing by $g_kg_{k-1}$, the
left side becomes
\[
  2(1+\eps_k)-(1+\eps_{k-1})-(1+\eps_{k+1})\,
  \frac{g_{k+1}g_{k-2}}{g_kg_{k-1}}
  \;=\;2(1+\eps_k)-(1+\eps_{k-1})-(1+\eps_{k+1})\,\frac{\rho_{k+1}}{\rho_{k-1}},
\]
and $\rho_{k+1}/\rho_{k-1}=\lambda_k$ by the closed form of $\rho$.
Collecting terms gives the identity.  For the reduction, write
$\lc_P(k)=\lc_F(k)+\lc_{xg}(k)+\mathrm{defect}(k)$ with
$\mathrm{defect}(k)=2F_kg_{k-1}-F_{k-1}g_k-F_{k+1}g_{k-2}$ (the mixed
terms of the square minus the two mixed products); $\lc_{xg}(k)\ge0$
because shifted binomial sequences are log-concave.  If
\eqref{eq:master} holds then by the identity
\[
  \mathrm{defect}(k)\;\ge\;-\,g_kg_{k-1}\,\rho_km_k(1+\eps_k)^2
  \;=\;-\,m_k\,g_k^2(1+\eps_k)^2\;=\;-\,m_kF_k^2\;\ge\;-\lc_F(k)
\]
by Theorem~\ref{thm:a2}, using $g_kg_{k-1}\rho_k=g_k^2$ and
$F_k=(1+\eps_k)g_k$.  Hence $\lc_P(k)\ge0$.
\end{proof}

\section{Per-class growth bounds}\label{sec:classes}

Group the classes by $(s,a)$ with $s=|S|$ and $a=\sum_{i\in S}c_i$:
the class mass is $w_{s,a}(k)=\nu_{s,a}\,u[k-s]/g_k$ with
$u=(1+x)^a(1+2x)^{C-a}$ and $\nu_{s,a}$ the number of subsets
realizing $(s,a)$, so that $W_s(k)=\sum_aw_{s,a}(k)$.  All bounds below concern the
\emph{growth} $\gamma(k)=w(k)/w(k-1)=r^u(k-s)/\rho_k$, where
$r^u(j)=u_j/u_{j-1}$.

\begin{lemma}[E1: ratio sandwich]\label{lem:e1}
For $u=(1+x)^a(1+2x)^{C-a}$ and all $j$ in the support,
\[
  \frac{C-j+1}{j}\;\le\;r^u(j)\;\le\;\frac{2(C-j+1)}{j}.
\]
\end{lemma}

\begin{proof}
$(1+x)\le_{\mathrm{lr}}(1+2x)$ (the coefficient ratio $1:2$ exceeds
$1:1$), and likelihood-ratio order between nonnegative sequences is
preserved by convolution with a common PF$_2$ (log-concave, no
internal zeros) factor \cite[Thm.~1.C.9]{ShakedShanthikumar};
equivalently, a two-line Cauchy--Binet check.  Applying this $a$
resp.\ $C-a$ times,
$(1+x)^C\le_{\mathrm{lr}}u\le_{\mathrm{lr}}(1+2x)^C$, and
$X\le_{\mathrm{lr}}Y$ implies $r^X(j)\le r^Y(j)$ wherever defined; the
two outer ratios are $(C-j+1)/j$ and $2(C-j+1)/j$.
\end{proof}

Two corollaries used throughout: $\eps_{k-1}\le2\eps_k$ on the band,
and $W_s(k)\le\binom hs q_k^s$ with $q_k=k/(2(C-k+1))$.

\begin{lemma}[Exact window means]\label{lem:means}
Write $u=(1+x)\tilde u$ with $\tilde u=(1+x)^{a-1}(1+2x)^{C-a}$, and
let $I$ (resp.\ $I'$) be the number of $(1+x)$-factors used at level
$j-1$ (resp.\ $j$) under the coefficient measure of $u$.  Then
\[
  \mathbb E[I]=\frac{a}{1+r^{\tilde u}(j-1)},\qquad
  \mathbb E[I']=\frac{a}{1+r^{\tilde u}(j)},
\]
and consequently, by Lemma~\ref{lem:e1} applied to $\tilde u$,
\[
  \mathbb E[I]\;\ge\;\frac{a(j-1)}{(j-1)+2(C-j+1)},\qquad
  \mathbb E[I']\;\le\;\frac{aj}{C}.
\]
\end{lemma}

\begin{proof}
$\sum_i i\binom ai v_{j-1-i}=a\sum_i\binom{a-1}{i-1}v_{j-1-i}
=a\,\tilde u_{j-2}$ for $v=(1+2x)^{C-a}$, while
$u_{j-1}=\tilde u_{j-1}+\tilde u_{j-2}$; divide.  The second identity
is the same computation one level up.
\end{proof}

\begin{lemma}[Jensen floor and ceiling]\label{lem:jensen}
With $r^v(m)=2(C-a-m+1)/m$,
\[
  r^u(j)\;\ge\;r^v\!\bigl(j-\mathbb E[I]\bigr)
  \qquad\text{and}\qquad
  r^u(j)\;\le\;r^v\!\bigl(j-\mathbb E[I']\bigr).
\]
Combined with Lemma~\ref{lem:means} these give closed-form rational
bounds $\mathrm{fl}(s,a,k)\le\gamma(k)\le\mathrm{ce}(s,a,k)$ for every
class growth.
\end{lemma}

\begin{proof}
Conditioning on $I$, $r^u(j)=\mathbb E_\omega[r^v(j-I)]$ where
$\omega$ is the level-$(j-1)$ window measure; $i\mapsto r^v(j-i)$ is
convex and increasing, so Jensen plus the floor on $\mathbb E[I]$ give
the first bound.  For the second, $1/r^u(j)=\mathbb
E_{\omega'}[1/r^v(j-I')]$ under the level-$j$ window measure, and
$i\mapsto 1/r^v(j-i)$ is convex and decreasing; Jensen plus the
ceiling on $\mathbb E[I']$ finish.
\end{proof}

\paragraph{The charge decomposition.}
The left side of \eqref{eq:master} splits exactly over classes:
\[
  \lambda_kd_k-d_{k-1}
  =\sum_{s,a}w_{s,a}(k)\Bigl[\lambda_k(\gamma_+-1)-\bigl(1-\tfrac1{\gamma_0}\bigr)\Bigr],
\]
where $\gamma_0,\gamma_+$ are the class growths into $k$ and $k+1$.
Each bracket admits two upper bounds (``strategies''):
\[
  \mathrm{cap}_A=\lambda_k(\mathrm{ce}_+-1)^++\frac{(1-\mathrm{fl})^+}{\mathrm{fl}},
  \qquad
  \mathrm{cap}_B=\lambda_k\Bigl(\mathrm{ce}_0\,\sigma_k
  +\frac{\mathrm{dev}^2}{\mathrm{fl}}\Bigr)
  +(1-\lambda_k)\frac{(1-\mathrm{fl})^+}{\mathrm{fl}},
\]
with $\sigma_k=\rho_k/\rho_{k+1}-1$ and
$\mathrm{dev}=\max(\mathrm{ce}_0-1,\,1-\mathrm{fl})$; strategy $B$
uses the second-difference identity
$\gamma_+-2+1/\gamma_0=(\gamma_+-\gamma_0)+(\gamma_0-1)^2/\gamma_0$
together with $\gamma_+\le\gamma_0(1+\sigma_k)$ (log-concavity of $u$).
Since both are upper bounds for the \emph{same} bracket, the
per-class minimum $\min(\mathrm{cap}_A,\mathrm{cap}_B)$ is a valid
charge, and
\begin{equation}\label{eq:charges}
  \lambda_kd_k-d_{k-1}\;\le\;\sum_sW_s(k)\,\mathrm{cap}(s),\qquad
  \mathrm{cap}(s)=\max_a\min(\mathrm{cap}_A,\mathrm{cap}_B).
\end{equation}

\section{The demotion flow}\label{sec:flow}

Fix a band position $k$ and define the \emph{demotion flow}
\[
  G(t)\;=\;\Bigl[\prod_{i=1}^h\bigl(E_i+t\,b_i\bigr)\Bigr]_k,
  \qquad \phi(t)=\ln G(t),\qquad t\in[0,2].
\]
Then $G(1)=F_k$, $G(0)=g_k$, and expanding the product over which
factors contribute $b_i$,
\[
  G(t)\;=\;g_k\sum_{s\ge0}W_s(k)\,t^s ,
\]
so the flow is the generating function of the class-mass profile.
Three exact facts:
\[
  \phi(1)-\phi(0)=\ln(1+\eps_k),\qquad
  \phi'(1)=\sum_{i=1}^h\frac{\delta_i}{1+\delta_i},\qquad
  \sum_s sW_s(k)=(1+\eps_k)\,\phi'(1),
\]
where $\delta_i=(F-F^{[i]})_k/F^{[i]}_k$ and $F^{[i]}$ is $F$ with hub
$i$ demoted ($H_i\to E_i$); the second uses
$b_i\prod_{l\ne i}H_l=F-F^{[i]}$.

\begin{proposition}[Flow facts]\label{prop:flow}
Suppose $\phi''\le0$ on $[0,2]$ at position $k$.  Then
\begin{align*}
  \textup{(P5)}\quad &\textstyle\sum_s sW_s\le(1+\eps)\ln(1+\eps),\\
  \textup{(m2)}\quad &\textstyle\sum_s s(s-1)W_s\le(1+\eps)\ln^2(1+\eps),\\
  \textup{(envelope +)}\quad
  &\textstyle\sum_{s\ge1}W_st^s\le(1+\eps)^t-1\qquad(1\le t\le2),\\
  \textup{(envelope --)}\quad
  &\textstyle\sum_{s\ge1}W_st^s\ge(1+\eps)^t-1\qquad(0\le t\le1).
\end{align*}
\end{proposition}

\begin{proof}
Concavity gives $\phi'(1)\le\phi(1)-\phi(0)=\ln(1+\eps)$, which is
(P5) after multiplying by $1+\eps$; $G''G\le(G')^2$ at $t=1$ gives
$\sum_ss(s-1)W_s\le(1+\eps)\phi'(1)^2$, which is (m2) by (P5).  For
the envelopes: on $[1,2]$, $\phi(t)\le\phi(1)+(t-1)\phi'(1)\le\ln
F_k+(t-1)\ln(1+\eps)$ (tangent line plus (P5)), and on $[0,1]$,
$\phi(t)\ge(1-t)\phi(0)+t\phi(1)$ (chord); exponentiate and divide by
$g_k$.\;
\end{proof}

\begin{proposition}[Diagonal dominance]\label{prop:dd}
With $A_l=E_l+tb_l$, $x_i=[b_i\prod_{l\ne i}A_l]_k$ and
$y_{ij}=[b_ib_j\prod_{l\ne i,j}A_l]_k$,
\[
  (G')^2-GG''=\Bigl(\sum_ix_i\Bigr)^2-G\sum_{i\ne j}y_{ij}
  =\sum_ix_i^2+\sum_{i\ne j}\bigl(x_ix_j-G\,y_{ij}\bigr),
\]
so hypothesis (ii) of Theorem~\ref{thm:reduction} is the aggregate
inequality $\sum_{i\ne j}(G\,y_{ij}-x_ix_j)\le\sum_ix_i^2$, and
$\phi''(t)\le0$ holds at any $(k,t)$ where all the $x_i$ are equal
and the symmetric form \eqref{eq:sym} holds for every pair: there
each ordered pair contributes at most $x^2/(h-1)$, and the
$h(h-1)$ ordered pairs sum to exactly $\sum_ix_i^2$.
\end{proposition}

\begin{proof}
The identity is the product rule for
$G'=\sum_ix_i$-type derivatives of the multiaffine flow; the
symmetric case is the stated count.
\end{proof}

\begin{remark}[Pairwise splits fail]\label{rem:nopairwise}
The natural per-pair sufficient conditions for the aggregate are
false in general.  Allocating the budget $\sum_ix_i^2$ uniformly,
$G\,y_{ij}-x_ix_j\le(x_i^2+x_j^2)/(2(h-1))$ (exact at symmetric
points), is violated by honest in-band instances at $1$-heavy
backgrounds --- smallest witness $S(1,27,1^{16})$, $n=107$, at
$k=5$; worst observed factor $7.5$ at $(1,128,1^{256})$, $k=10$ ---
while the aggregate holds there with the margin growing without
bound.  Any pairwise route to hypothesis (ii) for mixed multisets
must therefore allocate adaptively; the aggregate is the
invariant object.
\end{remark}

There is, however, a sufficient route that eliminates the flow
parameter entirely.  Write $G(t)=\sum_sW_s(k)\,t^s$, where
\[
  W_s(k)=\sum_{|S|=s}\Bigl[\,\prod_{i\in S}b_i\prod_{i\notin S}E_i\,
  \Bigr]_k
\]
is the class-mass profile (contiguous support $[0,\min(h,k)]$ at
band positions).

\begin{proposition}[Factorial criterion]\label{prop:ladder}
Let $G(t)=\sum_{s}W_st^s$ with $W_s\ge0$ and contiguous support.  If
the sequence $(s!\,W_s)$ is log-concave on the support, then $G$ is
log-concave on $(0,\infty)$.  Consequently hypothesis (ii) at
position $k$ follows from the \emph{ladder}
\begin{equation}\label{eq:ladder}
  s\,W_s(k)^2\;\ge\;(s+1)\,W_{s-1}(k)\,W_{s+1}(k),
  \qquad 1\le s\le\min(h,k)-1 .
\end{equation}
\end{proposition}

\begin{proof}
Log-concavity of $(s!W_s)$ says $\rho_s:=sW_s/W_{s-1}$ is
nonincreasing on the support.  Under the tilted law
$\mathbb P_t(S=s)=W_st^s/G(t)$,
\[
  (\ln G)'(t)=\frac{G'(t)}{G(t)}
  =\sum_s\frac{(s+1)W_{s+1}t^s}{G(t)}
  =\mathbb E_t\bigl[\rho_{S+1}\bigr]
\]
(with $\rho:=0$ above the support).  For $t'>t$ the likelihood ratio
$\mathbb P_{t'}(s)/\mathbb P_t(s)\propto(t'/t)^s$ is increasing in
$s$, so $\mathbb P_{t'}$ dominates $\mathbb P_t$ in the
likelihood-ratio order, hence stochastically
\cite{ShakedShanthikumar}; since $\rho_{S+1}$ is nonincreasing,
$\mathbb E_t[\rho_{S+1}]$ is nonincreasing in $t$.  Thus $(\ln G)'$
is nonincreasing on $(0,\infty)$.
\end{proof}

The ladder \eqref{eq:ladder} is tight at its base and unifies the
results of this paper:
the rung $s=1$ reads $W_1^2\ge2W_0W_2$, which is \emph{identically}
hypothesis (ii) at $t=0$ (so the bottom rung is necessary); at $h=2$
the ladder is that single rung, which is the inequality proved in
Theorem~\ref{thm:h2mixed}; and at uniform multisets, where
$W_s=\binom hsu_s$ with $u_s=[b^sE^{h-s}]_k$, rung $s$ becomes
$u_s^2\ge u_{s-1}u_{s+1}(h-s)/(h-s+1)$, whose first case is the
symmetric E-core \eqref{eq:sym}.  Each bracket $u_s$ is a polynomial
in $c$ at fixed $(k,h,s)$, so the higher rungs are certifiable by
the method of Section~\ref{sec:polyc}.  Exact scans
(Section~\ref{sec:verification}) found no rung failure anywhere ---
and located the binding rung at $s=1$ throughout --- so the ladder
is the route we pursue for mixed multisets
(Section~\ref{sec:open}).  We also record the polynomial identity
behind rung $1$: with $\Phi_s=\sum_{|S|=s}\prod_{i\in
S}b_i\prod_{i\notin S}E_i$ (so $W_s=[\Phi_s]_k$), the cross terms in
$\Phi_1^2$ cancel against $2\Phi_0\Phi_2$ exactly, leaving
\[
  \Phi_1^2-2\Phi_0\Phi_2
  =\sum_ib_i^2\Bigl(\prod_{l\ne i}E_l\Bigr)^{\!2},
\]
which is coefficientwise nonnegative; rung $1$ is the localization
of this identity to the $(k,k)$ product of coefficients.

The class masses also carry an exact majorization structure.

\begin{proposition}[Schur-convexity of the class masses]
\label{prop:wschur}
At fixed $(C,h,k,s)$, the class mass $W_s$ is Schur-convex on
multisets with sum $C$ and $h$ parts; in particular
$W_s(c)\ge W_s(\mathrm{bal}(C,h))$ for every multiset $c$, where
$\mathrm{bal}(C,h)$ is the balanced two-value partition.
\end{proposition}

\begin{proof}
A Robin Hood transfer toward balance on slots $(i,j)$ with
$c_i<c_j$ induces, on the multiset of $s$-subset sums
$\sigma_S=\sum_{l\in S}c_l$, the simultaneous transfers
$(\sigma_{S\cup i},\sigma_{S\cup j})\to(\sigma_{S\cup
i}{+}1,\sigma_{S\cup j}{-}1)$ over all $(s{-}1)$-subsets $S$
avoiding $i,j$ --- each toward balance, since
$\sigma_S+c_i<\sigma_S+c_j$; subsets containing both or neither
are unchanged.  Hence the subset-sum multiset on the spread side
majorizes the balanced side's.  The summand
$u(\sigma)=[x^{k-s}](1+x)^\sigma(1+2x)^{C-\sigma}$ is convex in
$\sigma$, by the exact second difference
$\Delta^2u=[x^{k-s-2}](1+x)^\sigma(1+2x)^{C-\sigma-2}\ge0$ (linear
at $k{-}s{=}1$, constant at $k{-}s{=}0$), and Karamata's inequality
concludes.
\end{proof}

Thus both sides of every rung grow as the multiset spreads, and
the open content of the ladder for mixed multisets is a
relative-rate statement (the domination lemma of
Section~\ref{sec:open}).

\section{The E-core and the flow concavity at
\texorpdfstring{$h=2$}{h=2}}\label{sec:p6}

The two reduction lemmas of this section concern the symmetric form
\eqref{eq:sym}; they power the uniform-case certification, where by
Proposition~\ref{prop:dd} the symmetric form at every pair is
equivalent to the flow concavity.  For mixed multisets no pairwise
reduction is available (Remark~\ref{rem:nopairwise}) and the
aggregate diagonal $t$-family must be treated as such
(Section~\ref{sec:open}); the exception is $h=2$, where the
aggregate \emph{is} the single-pair inequality and is proved below
for even total degree.

\begin{lemma}[Vertex reduction]\label{lem:vertex}
The supremum over $(t_1,\dots,t_h)\in[0,2]^h$ of the symmetric-form
defect $[b_ib_jW][A_iA_jW]/\bigl([b_iA_jW][b_jA_iW]\bigr)$ with
$A_l=E_l+t_lb_l$ is attained with each $A_l\in\{E_l,\,E_l+2b_l\}$.
\end{lemma}

\begin{proof}
The defect is a ratio of polynomials with positive coefficients,
multiaffine in the $t_l$; as a function of each $t_l$ separately it is
a M\"obius transformation with no pole on $[0,\infty)$, hence monotone
there, so the supremum over the box is attained at a vertex.
\end{proof}

\begin{lemma}[Slot expansion]\label{lem:slot}
For any nonnegative background $W$, any $d\ge0$, and $H_2:=E+2b$, the
inequalities with $X_i,X_j\in\{E,H_2\}$ all follow termwise from the
single core case $X_i=X_j=E$ with the same $W$ and constant $1+d$.
\end{lemma}

\begin{proof}
Expand $[H_2X]=[EX]+2[bX]$ in each slot; every cross term appears on
both sides with matching or dominating coefficients; the only
nontrivial comparison is the core case.
\end{proof}

\begin{proof}[Proof of Theorem~\ref{thm:h2}]
By Lemmas~\ref{lem:vertex} and~\ref{lem:slot} it suffices to prove
$[b^2]_k[E^2]_k\le 2([bE]_k)^2$.  All three brackets are explicit:
$[b^2]_k=\binom{2c}{k-2}$, $[E^2]_k=\binom{2c}k2^k$, and since
$(1+x)^c(1+2x)^c=(1+3x+2x^2)^c$,
\[
  [bE]_k\;=\;T_{k-1},\qquad
  T_n:=[x^n](1+3x+2x^2)^c=\sum_m\binom cm\binom c{n-m}2^{\,n-m}.
\]
By the symmetry $m\leftrightarrow n-m$ of $\binom cm\binom c{n-m}$ and
AM--GM, followed by Vandermonde,
\[
  T_n=\tfrac12\sum_m\binom cm\binom c{n-m}\bigl(2^m+2^{\,n-m}\bigr)
  \;\ge\;2^{n/2}\sum_m\binom cm\binom c{n-m}
  \;=\;2^{n/2}\binom{2c}n .
\]
Hence, using log-concavity of the binomial row $\binom{2c}\cdot$,
\[
  2\,T_{k-1}^2\;\ge\;2^{k}\binom{2c}{k-1}^2
  \;\ge\;2^{k}\binom{2c}{k-2}\binom{2c}k\;=\;[b^2]_k\,[E^2]_k. \qedhere
\]
\end{proof}

\begin{lemma}[$t$-cancellation]\label{lem:tcancel}
Fix any background polynomial $W$ and a position $k$, and write
$A_i=E_i+t\,b_i$, $A_j=E_j+t\,b_j$ on the two pair slots.  Then
\[
  [b_ib_jW]_k[A_iA_jW]_k-[b_iA_jW]_k[b_jA_iW]_k
  \;=\;[b_ib_jW]_k[E_iE_jW]_k-[b_iE_jW]_k[b_jE_iW]_k ,
\]
independent of $t$.  Consequently, since the brackets
$[b_iA_jW]_k=[b_iE_jW]_k+t[b_ib_jW]_k$ are nondecreasing in $t$, any
upper bound on the pairwise defect that is monotone in the
cross-brackets need only be verified at $t=0$ on the pair slots.
\end{lemma}

\begin{proof}
Brackets are linear in each slot polynomial, so with
$P=[b_ib_jW]_k$, $Q_0=[E_iE_jW]_k$, $X_0=[b_iE_jW]_k$,
$Y_0=[b_jE_iW]_k$, the four $t$-dependent brackets are
$Q_0+t(X_0+Y_0)+t^2P$, $X_0+tP$, and $Y_0+tP$; expanding,
$P\bigl(Q_0+t(X_0+Y_0)+t^2P\bigr)-(X_0+tP)(Y_0+tP)=PQ_0-X_0Y_0$.
Equivalently, the matrix
$M(t)=\bigl(\begin{smallmatrix}P&Y\\X&Q\end{smallmatrix}\bigr)$
satisfies $M(t)=(I+te_2e_1^{\mathsf T})\,M(0)\,(I+te_1e_2^{\mathsf T})$,
a unipotent congruence, which preserves the determinant.
\end{proof}

\begin{lemma}[Schur minimality of the balanced
split]\label{lem:schurS}
For $a+b=s$ let $S(a,b)=(1+x)^a(1+2x)^b+(1+x)^b(1+2x)^a$.  Then for
$a<b$,
\[
  S(a,b)-S(a+1,b-1)
  =x\,\bigl[(1+x)(1+2x)\bigr]^a
  \Bigl[(1+2x)^{\,b-a-1}-(1+x)^{\,b-a-1}\Bigr]\;\ge\;0
\]
coefficientwise; hence the coefficients of $S$ are minimized, among
the splits of $s$, at the balanced one.
\end{lemma}

\begin{proof}
With $d=b-a$, $S(a,b)=[(1+x)(1+2x)]^a\,T_d$ where
$T_d=(1+x)^d+(1+2x)^d$, and
$T_d-(1+x)(1+2x)T_{d-2}
=(1+x)^{d-1}\bigl[(1+x)-(1+2x)\bigr]
+(1+2x)^{d-1}\bigl[(1+2x)-(1+x)\bigr]
=x\bigl[(1+2x)^{d-1}-(1+x)^{d-1}\bigr]$,
whose coefficients are nonnegative because $(1+2x)^{d-1}$ dominates
$(1+x)^{d-1}$ coefficientwise.
\end{proof}

\begin{lemma}[Odd core]\label{lem:oddcore}
Let $s=2m+1$, $m\ge1$, and $T_n=[x^n](1+3x+2x^2)^m$.  Then for all
$k\ge2$,
\[
  2\binom s{k-2}2^k\binom sk\;\le\;\bigl(2T_{k-1}+3T_{k-2}\bigr)^2 .
\]
\end{lemma}

\begin{proof}
For $k>s$ the left side vanishes; let $2\le k\le s$ and put
$a=\binom{2m}{k-2}\ge1$, $b=\binom{2m}{k-1}\ge1$.  By Pascal's rule
and log-concavity of the row $\binom{2m}\cdot$,
\[
  \binom s{k-2}=a+\binom{2m}{k-3}\le a\Bigl(1+\frac ab\Bigr),\qquad
  \binom sk=b+\binom{2m}k\le b\Bigl(1+\frac ba\Bigr),
\]
so $\binom s{k-2}\binom sk\le(a+b)^2$.  The AM--GM/Vandermonde bound
from the proof of Theorem~\ref{thm:h2}, $T_n\ge2^{n/2}\binom{2m}n$,
gives
\[
  2T_{k-1}+3T_{k-2}\;\ge\;2^{k/2}\Bigl(\sqrt2\,b+\tfrac32\,a\Bigr)
  \;\ge\;2^{k/2}\sqrt2\,(a+b),
\]
the last step because $\tfrac32>\sqrt2$.  Squaring,
$(2T_{k-1}+3T_{k-2})^2\ge2^{k+1}(a+b)^2
\ge2\binom s{k-2}2^k\binom sk$.
\end{proof}

The constant $3$, inherited from $T_1=2+3x$ at the balanced split, is
exactly what the proof needs: a cross term $\beta x$ with
$\beta/2<\sqrt2$ would fail.  The lemma is also exactly verified for
all $s\le161$ and all $k$, with worst ratio $0.5217$ --- the same
strength as the even core ($0.5218$).

\begin{proof}[Proof of Theorem~\ref{thm:h2mixed}]
At $h=2$ there is no background and
$(G')^2-GG''=x_1^2+x_2^2+2(x_1x_2-Gy)$ with $G=[A_1A_2]_k$,
$y=[b_1b_2]_k$, $x_i=[b_iA_j]_k$, so \eqref{eq:flowcc} reads
$2(Gy-x_1x_2)\le x_1^2+x_2^2$.  By Lemma~\ref{lem:tcancel} (with
$W=1$) the left side is independent of $t$ while the right side is
nondecreasing in $t$, so it suffices to prove, at $t=0$,
\[
  2\,[b_1b_2]_k\,[E_1E_2]_k\;\le\;(X_0+Y_0)^2,
  \qquad X_0=[b_1E_2]_k,\ Y_0=[b_2E_1]_k .
\]
Here $[b_1b_2]_k=\binom s{k-2}$ and $[E_1E_2]_k=2^k\binom sk$ depend
only on $s=c_1+c_2$, while $X_0+Y_0=[x^{k-1}]\,S(c_1,c_2)$ is
minimized at the balanced split by Lemma~\ref{lem:schurS}.  For even
$s$,
$[x^{k-1}]S(s/2,s/2)=2\,[x^{k-1}](1+3x+2x^2)^{s/2}=2\,T_{k-1}$, and
the display in the proof of Theorem~\ref{thm:h2}, at $c=s/2$, states
exactly $\binom s{k-2}2^k\binom sk\le2\,T_{k-1}^2$, i.e.
$2\,[b_1b_2]_k[E_1E_2]_k\le(2T_{k-1})^2\le(X_0+Y_0)^2$.  For odd
$s$, $S\bigl(\tfrac{s-1}2,\tfrac{s+1}2\bigr)
=(1+3x+2x^2)^{(s-1)/2}\bigl((1+x)+(1+2x)\bigr)$, so
$[x^{k-1}]S_{\mathrm{bal}}=2T_{k-1}+3T_{k-2}$ with $T$ at exponent
$(s-1)/2$, and Lemma~\ref{lem:oddcore} closes the chain.
\end{proof}

\begin{proof}[Proof of Theorem~\ref{thm:limit}]
With $A=2\ell+2$, $B=2\ell+3$, the bounds $\ln(1-x)\le-x$ and
$\ln(1+x)\le x$ give
$\ln(1+\Delta_\infty)\le-\tfrac1k+\tfrac2A-\tfrac k{B^2}$, and by
AM--GM $\tfrac1k+\tfrac k{B^2}\ge\tfrac2B$, so
$\ln(1+\Delta_\infty)\le 2/(AB)\le 1/6$.  Since $e^x-1\le1.1x$ on
$(0,1/6]$,
\[
  (h-1)\,\Delta_\infty\;\le\;(\ell+1)\cdot\frac{2.2}{(2\ell+2)(2\ell+3)}
  \;=\;\frac{1.1}{2\ell+3}\;\le\;0.22 .\qedhere
\]
\end{proof}

\begin{proof}[Derivation of $\Delta_\infty$]
In the uniform all-$E$ background, fix $k$ and let $c\to\infty$.  The
needed bracket types are $[b^2E^{h-2}]_k$, $[E^h]_k$, and
$[bE^{h-1}]_k$.  In each finite convolution sum,
\[
  \binom{ac}m2^m=\frac{(2ac)^m}{m!}\bigl(1+O(m^2/(ac))\bigr),
\]
and therefore
\begin{align*}
  [E^h]_k
  &=\binom{hc}k2^k\sim\frac{(2hc)^k}{k!},\\
  [bE^{h-1}]_k
  &=\sum_m\binom cm\binom{(h-1)c}{k-1-m}2^{k-1-m}\\
  &\sim\frac{(2(h-1)c)^{k-1}}{(k-1)!}
  \sum_m\binom{k-1}m\Bigl(\frac1{2(h-1)}\Bigr)^m,
\end{align*}
and the binomial sum equals $\bigl(1+\tfrac1{2(h-1)}\bigr)^{k-1}$,
giving $[bE^{h-1}]_k\sim\bigl((2h-1)c\bigr)^{k-1}/(k-1)!$; similarly
$[b^2E^{h-2}]_k\sim\bigl(2(h-1)c\bigr)^{k-2}/(k-2)!$ via
$2(h-2)\bigl(1+\tfrac1{h-2}\bigr)=2(h-1)$.  Hence, with $B=2h-1$,
\[
  \Delta_\infty+1
  =\frac{(2(h-1))^{k-2}(2h)^k}{B^{2k-2}}\cdot\frac{(k-1)!^2}{(k-2)!\,k!}
  =\Bigl(1-\frac1{B^2}\Bigr)^{k-2}\Bigl(1+\frac1B\Bigr)^2
  \Bigl(1-\frac1k\Bigr),
\]
and the identity $(1-B^{-2})^{-2}(1+B^{-1})^2=(1+\tfrac1{B-1})^2$
converts this to the form displayed in Theorem~\ref{thm:limit} with
$B-1=2\ell+2$.
\end{proof}

The finite-$c$ statement of Theorem~\ref{thm:cert} is proved by the
polynomial certification method of Section~\ref{sec:verification}.

\section{Certificates and machine verification}\label{sec:verification}

This section describes the machine-verified components and the
protocol under which they were produced and audited.  All
verifications are exact: every inequality is decided in rational
arithmetic (Python \texttt{fractions.Fraction} or \texttt{sympy}
integer polynomials), with no floating point in any verified
statement.

\subsection{Certificates for the master inequality}

For each multiset and band position, the master inequality
\eqref{eq:master} is established through \eqref{eq:charges} by one of
two explicit certificates.

\emph{Certificate T1} (unconditional).  From the identity
$1+\eps_{k+1}=(r_{k+1}/\rho_{k+1})(1+\eps_k)$ and $r_{k+1}\le r_k$
(log-concavity of $F$, by Lemmas~\ref{lem:a0} and~\ref{lem:a1}),
\[
  \lambda_k(1+\eps_k)\Bigl[(1+\sigma_k)\frac{1+\eps_k}{1+\eps_{k-1}}-1\Bigr]
  +(\eps_{k-1}-\eps_k)^+\;\le\;\text{RHS of \eqref{eq:master}}
\]
suffices, a single rational inequality depending only on
Lemmas~\ref{lem:a0}, \ref{lem:a1} and Theorem~\ref{thm:a2}.

\emph{Certificate T2} (uses the flow facts, hence the flow
concavity \eqref{eq:flowcc} at the same position).  A quadratic
majorant
$\mu+y_1s+y_2s(s-1)\ge\mathrm{cap}(s)$, anchored at $s=1,2$ with $y_2$
the maximal second divided difference (clamped so $y_1\ge0$), is
played against the \emph{rationalized} moment rows
\[
  \sum_ssW_s\le R_1=\frac{\eps(2+\eps)}{2},\qquad
  \sum_ss(s-1)W_s\le R_2=\frac{\eps^2(2+\eps)^2}{4(1+\eps)},
\]
which follow from Proposition~\ref{prop:flow} and the trapezoid bound
$\ln(1+x)\le x(2+x)/(2(1+x))$; the residual tail is covered by the
class bounds $W_s\le\binom hsq_k^s$ or the $t=2$ envelope, whichever
is cheaper.  Pointwise domination is checked exactly for every $s$,
and the certificate obligation is one rational comparison.

These certificate families were located by linear programming
(maximize the charge sum over all profiles $(W_s)$ satisfying the
proved constraints; the optimum is $25.97\%$ below the requirement at
the worst point of the grid), and the dual solutions were then
replaced by the closed-form constructions above.  Verification: the
grid $\mathcal G$ consists of $305$ distinct multisets ($215$ uniform,
covering every $(c,h)$ with $c\le8$, $h(1+c)\le120$ or $c\le10$,
$n\le160$; and $90$ random mixed multisets with $c_i\le12$); at every
band position of every member, T1 or T2 passes in exact arithmetic;
worst certificate ratio $0.4352$.

The separate Band~B verifier proves the finite residual left by
Section~\ref{sec:bandB}: after the explicit $C\ge55$ threshold, it
enumerates every leaf-heavy multiset with $C\le54$ and nonempty
Band~B.  This is $24{,}394$ multisets and $54{,}049$ coefficient
positions, all checked directly in exact integer arithmetic with zero
failures.  The same script records the large-threshold residuals through
$C=1000$ and their positive boundary minimum at $(C,h)=(55,35)$; the
all-$C$ part is the displayed quotient monotonicity check in
Section~\ref{sec:bandB}.

\subsection{Polynomial certification of the E-core}\label{sec:polyc}

At fixed $(k,h)$ and background pattern $W=E^{h-2-w}(E+2b)^w$ in the
uniform case, each bracket of \eqref{eq:sym} is a polynomial in $c$
of degree at most $k$ (each factor coefficient is a polynomial in $c$;
the bracket is a finite convolution).  The claim ``\eqref{eq:sym}
holds for \emph{all integers} $c\ge1$'' is therefore decidable
exactly: write $R(c)=h\,[bEW]_k^2-(h-1)[b^2W]_k[E^2W]_k$, check the
leading coefficient is positive, choose $B$ at least the largest real
root (any upper bound suffices; we use a numerically located bound,
then \emph{prove} positivity beyond it by verifying that all
coefficients of the Taylor shift $R(c+B)$ are nonnegative), and
evaluate $R$ exactly at every integer in $[1,B]$.  The numeric root
location plays no role in the rigor; it only sizes the verified
range.  This certified the $996$ instances of Theorem~\ref{thm:cert}
with zero failures.

\subsection{Audit}\label{sec:audit}

Because nearly every verification script was written by the same
author as the claim it checks, the full chain was subjected to an
adversarial audit: independent re-derivation of every load-bearing
identity; an exact six-way sub-bound audit of the charge machinery
over $10{,}751$ class instances; independent reconstruction of the
certification brackets at random points; and end-to-end brute-force
unimodality of every tree in the claimed scope by the original
dynamic-programming code.  The audit found two implementation-level
flaws (an incorrectly indexed ceiling inside $\mathrm{dev}$, masked by
unrelated slack, and a finite sentinel acting as a false bound at
class-birth positions), both repaired with all headline figures
unchanged, and no mathematical flaw.  The audit report and all scripts
are in the accompanying repository.

The same adversarial methodology was applied to the working
conjectures themselves, and twice replaced them.  The symmetric form
\eqref{eq:sym}, formulated from uniform data, was refuted at lopsided
pairs $(1,c,1^{h-2})$; its exact-allocation weighted replacement was
in turn refuted at $1$-heavy backgrounds $(c_1,c_2,1^A)$ --- a corner
the first sweep could not reach because it capped $h$ at $8$ ---
with smallest witness $S(1,27,1^{16})$, $n=107$, $k=5$.  The
surviving hypothesis is the aggregate \eqref{eq:flowcc} itself,
scanned in exact arithmetic across the flow parameter
$t\in\{0,\tfrac12,1,\tfrac32,2\}$: $14{,}410$ evaluations over the
families $(c,c,1^A)$, $(1,c,1^A)$, $(2,c,1^A)$ ($c\le128$,
$A\le256$), all $200$ pairwise-violation witnesses, and $40$ random
mixed multisets, with zero violations, worst ratio $+0.085$, and the
margin improving to $-5.3$ in probes to $A=1024$.  At uniform
multisets the aggregate ratio was confirmed to coincide with the
symmetric-form ratio (zero mismatches), which is the formal statement
that Theorems~\ref{thm:h2}, \ref{thm:limit} and~\ref{thm:cert}
certify instances of hypothesis (ii).

The ladder \eqref{eq:ladder} was scanned with the same rigor:
$5{,}298$ class-mass profiles in exact arithmetic --- uniform
multisets ($c\le12$, $h\le12$), the $(a,b,1^A)$ families including
every pairwise-violation witness ($c\le128$, $A\le256$), and $60$
random mixed multisets --- with \emph{zero} rung failures; worst
rung ratios $0.917$ (uniform), $0.971$ ($(128,128,1^{256})$,
$k=40$), $0.884$ (random mixed), always at the necessary rung
$s=1$.  A Schur analysis over \emph{all} partitions of $C$ into $h$ parts
shows the reduction landscape is a threshold phenomenon, not a
monotonicity: transfer-by-transfer monotonicity of the rung ratio
is \emph{false} ($11{,}278$ violations among $49{,}941$ exact
Robin Hood steps, margins to $-0.23$), and in $212$ of $1{,}082$
exhaustive $(C,h,k,s)$ cases the maximizing partition is not the
most balanced one.  A later sweep of the bulk-plus-giant family
$(a^{h-1},M)$ ($4{,}720$ instances to $h=100$, $C=1200$) found it
beating the balanced companion even deep in the tight regime
(balanced ratios up to $0.956$ beaten), so no domination statement
holds; what holds in every observation, without exception, is the
ladder itself --- the sweep maximum is $0.9865<1$, at the minimal
gap-two giant $(6^{99},8)$, $k=100$.

Moreover, the polynomial certification extends to the rungs on the
two-adjacent-value families $(a^p,(a{+}1)^q)$ --- the target of the
Ladder--Schur reduction: at fixed $(k,h,p,s)$ each class mass is a
polynomial in $a$ of degree at most $k-s$, and the rung was
certified for \emph{all} integers $a\ge1$ at every instance with
$3\le h\le10$, $1\le p\le h$, $k\le28$: $7{,}866$ instances, zero
failures (with $p=h$ covering the uniform higher rungs) --- the
ladder holds for every two-adjacent-value multiset within the reach
of the uniform certification grid.
As a consistency datapoint, the balanced partition $(1^4,2^{254})$
of $(C,h)=(512,258)$ at $k=40$ has rung-$1$ ratio $0.9729$,
slightly above the $0.9711$ of the lopsided $(128,128,1^{256})$ ---
the balanced partition is the worse one, as the reduction predicts.

\section{Assembly: proofs of Theorems~\ref{thm:scope}
and~\ref{thm:reduction}}\label{sec:assembly}

\begin{proof}[Proof of Theorem~\ref{thm:reduction}]
Fix the multiset and let $k$ be a band position \eqref{eq:band}.
Hypothesis (ii) is precisely $\phi''\le0$ on $[0,2]$ at $k$, hence
the flow facts hold (Proposition~\ref{prop:flow}), hence the validity
of certificate T2's
constraint rows; certificate T1 is valid unconditionally.  By the
verified certificates, \eqref{eq:master} holds at $k$; by
Lemma~\ref{lem:epsid}, $P$ is log-concave at $k$.  Log-concavity at
$k=1$ is universal (Section~\ref{sec:skeleton}), and band B is covered
by Section~\ref{sec:bandB}.  Thus $P$ is log-concave on
$[1,k_{\mathrm{dec}}-1]$ and nonincreasing on $[k_{\mathrm{dec}},D]$
by Levit--Mandrescu/Basit--Galvin, so Lemma~\ref{lem:seam} applies.
Hypothesis (i) enters exactly once, in the penultimate sentence: the
master inequality at $k$ is established by the verified certificate at
$(c_1,\dots,c_h;k)$.
\end{proof}

\begin{proof}[Proof of Theorem~\ref{thm:scope}]
$h=1$ is Lemma~\ref{lem:d1}.
For $h=2$: Theorem~\ref{thm:h2} supplies the symmetric form at every
position with constant $2=1+\tfrac1{h-1}$, hence hypothesis (ii) by
Proposition~\ref{prop:dd} (a single pair; the AM--GM step), so
Theorem~\ref{thm:reduction} applies to every $S(c,c)$ in the
certificate grid; for the eleven mixed members, hypothesis (ii) is
Theorem~\ref{thm:h2mixed} directly (band B is empty at $h=2$:
$k_{\mathrm{dec}}\le\lceil(2C+3)/3\rceil\le C$ for $C\ge3$).  For
the sixteen three-hub adjacent-two-value families, hypothesis (ii)
is supplied by the ladder certificates on $(a^p,(a{+}1)^q)$
($h=3$, all integer $a$, $k\le28\ge k_A$) with
Proposition~\ref{prop:ladder}, hypothesis (i) was verified per
multiset (worst certificate ratio $0.38$), and band B is empty
except at $(1,1,2)$, whose single position is checked directly.  For $3\le h\le10$: Theorem~\ref{thm:cert} supplies
the symmetric form, over the whole flow box by
Lemmas~\ref{lem:vertex} and~\ref{lem:slot}, at every band position
provided $k_A$ does not exceed the certified $k$-range, which holds
exactly for the table of Theorem~\ref{thm:scope}; at uniform
multisets the diagonal is symmetric, so Proposition~\ref{prop:dd}
yields hypothesis (ii) and Theorem~\ref{thm:reduction} applies.
\end{proof}

\section{Extremal results and the counterexample
search}\label{sec:extremal}

\begin{definition}[Balance]\label{def:balance}
For an interior index $k$ of a positive sequence $(a_k)$, the (local)
\emph{balance} is $\min(a_{k-1},a_{k+1})/a_k$.  A positive sequence
fails unimodality if and only if some index $k$ satisfies the global
condition $\min\bigl(\max_{i<k}a_i,\;\max_{i>k}a_i\bigr)>a_k$; local
balance $>1$ is the width-one special case, and every log-concavity
failure observed in trees to date has width exactly one, so local
balance is the operative diagnostic.
\end{definition}

A valley at $k$ is precisely an interior index with balance $>1$, so
searching for a unimodality counterexample is the optimization problem
``maximize balance over trees''; every valley bottom is in particular
a log-concavity failure.  Three structural facts organize the search
space.

\paragraph{The offset law.}
For uniform hub spiders, $I(S(c^h))=H_c^h+xE_c^h$, and over the range
$h\le25$, $c\le12$ ($n$ up to $256$) the unique log-concavity failure
sits at offset exactly $h-2$ below the top degree for every $c\ge4$,
always with width one.  The failure is a \emph{junction}: the last
index where the fat $xE^h$ hump still contributes against the thin
$H^h$ tail.  Balance peaks at $0.0439$ for $S(4^5)$ ($n=46$) and
decays in both directions.

\paragraph{A high-fragility tree.}
A beam search over junction trees built from all $2{,}975$ rooted
gadgets on at most $11$ vertices, seeded with hub-spider multisets,
produced the largest balance in this search corpus: a $48$-vertex tree
(root, four
$H_4$ hub gadgets, one $11$-vertex gadget) with a single
log-concavity failure of balance $0.04512$ --- more than double the
previous record ($0.0208$, for the family $T^*$) --- and \emph{every}
gadget extension beyond it strictly degrades the balance until the dip
disappears.  Its sequence is non-monic (leading coefficient $2$),
which did not occur among the cited Kadrawi--Levit examples.

\paragraph{Why the mechanism self-defends.}
Writing any tree as root plus rooted gadgets,
$P=\prod F_i+x\prod G_i$, a junction valley needs the $\prod F$ side
still rising where the $\prod G$ hump ends, while near the junction
the coefficients of $\prod F$ organize as elementary symmetric
functions $e_j$ of per-gadget weights.  Dip existence forces
$e_2<1$ while the achievable balance behaves like $e_3/e_2$, and
Newton--Maclaurin gives $e_3/e_2\lesssim\tfrac12\sqrt{e_2}$: the same
log-concavity phenomenon the conjecture asserts is what caps the only
known failure mechanism, empirically at balance ${\sim}0.05$ against
the required $>1$.  Moreover the decreasing-zone theorems protect the junction's
location: the observed failures sit at offset $h-2$ from the top
degree, while the Basit--Galvin decreasing zone covers the top
$D-\ell\approx D^2/(D+n)$ positions, and
$(h+C)^2>(h-2)(3C+h+1)$ for all $c_i\ge1$, so every position the
known mechanism can reach lies inside the zone where a valley
(an ascent $i_k<i_{k+1}$) is already impossible.  A unimodality
counterexample by this mechanism would need the dip pushed below the
decreasing zone, and the balance data above shows the mechanism
weakening, not strengthening, in that direction.

Two further observations close the accessible routes: nesting a
dipped tree as a gadget under a new root erases its failure (the
second-level $F+xG$ smoothing), and products of tree polynomials
smooth width-one cliffs away --- in exact scans through total size
$64$, every forest product checked was log-concave outright, and the
flat-tail/steep-cliff combination a non-unimodal product requires
never co-occurs in real tree polynomials.  All searches and caps are,
of course, evidence about accessible scales only; they are recorded
here because the structural caps (the offset law, the Maclaurin
tension, the decreasing-zone protection) are proof-shaped statements
that delimit where any counterexample must live: width $\ge2$
log-convex windows, for which no generator is known.

\section{Open problems}\label{sec:open}

\begin{conjecture}[Flow concavity]\label{conj:ecore}
The flow concavity \eqref{eq:flowcc} holds for all $h\ge2$, all
$c_1,\dots,c_h\ge1$, all band positions $k$, and all $t\in[0,2]$:
the demotion flow $G(t)=[\prod_l(E_l+tb_l)]_k$ is log-concave in $t$.
\end{conjecture}

By Theorem~\ref{thm:reduction} the conjecture implies unimodality of
$I(S)$ for every hub spider.  The evidence: Theorems~\ref{thm:h2},
\ref{thm:h2mixed}, \ref{thm:limit} and~\ref{thm:cert} (which prove it
at $h=2$ and at uniform multisets over the certified range), plus the
exact aggregate scans of Section~\ref{sec:verification} ($14{,}410$
instances, zero violations, worst ratio $0.085$).  We emphasize the
history recorded in Remark~\ref{rem:nopairwise}: two successive
pairwise strengthenings of this conjecture were refuted by
systematic probing, at lopsided pairs and at $1$-heavy backgrounds
respectively, while the aggregate form survives both corners with
growing margin.  The conjecture as stated is exactly what the
pipeline uses --- nothing weaker suffices for
Proposition~\ref{prop:flow}, and nothing stronger is true among the
natural candidates.

\begin{conjecture}[All-parameters certificates]\label{conj:g2}
The certificate inequalities T1/T2 of
Section~\ref{sec:verification} hold for every multiset
$(c_1,\dots,c_h)$ and every band position, not only on the verified
grid.  (All certificate ratios are observed monotone or saturating in
$(h,c,k)$, with worst ratio $0.4352$.)
\end{conjecture}

\begin{question}
Does the demotion-flow method extend to root-gadgets of independence
number $\ge3$ (the \emph{gadget dichotomy}), and thence to broader
classes of trees in the direction of the full
Alavi--Malde--Schwenk--Erd\H{o}s conjecture?
\end{question}

Beyond the two conjectures, three concrete work items would complete
the program for all hub spiders.  First,
Conjecture~\ref{conj:ecore} for mixed multisets with $h\ge3$, for
which Proposition~\ref{prop:ladder} supplies a $t$-free route: it
suffices to prove the ladder \eqref{eq:ladder}.  The supporting
structure established so far: the rungs are certified for
\emph{all} two-adjacent-value multisets $(a^p,(a{+}1)^q)$ in the
range $h\le10$, $k\le28$ (all integer $a$, Section
\ref{sec:verification}), and the exact scans show a threshold
landscape --- whenever a lopsided partition beats the balanced one,
all ratios in that case sit below $0.544$, while the
near-$1$ corners are dominated by balanced partitions.  Three structural shortcuts are provably unavailable: plain Schur
reduction (transfer monotonicity toward balance fails), its
conditional form (violating transfers exist at ratio $0.868$), and
a value-window bound (one outlier in a large balanced bulk leaves
the ratio at $0.944$).  What survives every exact probe
($\sim$25{,}000 cases) is endpoint domination in its canonical
conditioning: whenever the \emph{balanced companion's} rung ratio
is at least $0.7$, no multiset with the same $(C,h)$ beats it ---
the balanced companion is never beaten above ratio $0.66$ in any
scan, while at low ratios beats are common and harmless.  This
\emph{domination lemma} is the open core.  Two pieces of supporting
structure are in place: Proposition~\ref{prop:wschur} (both sides
of every rung grow away from balance, so domination is a
relative-rate statement), and the \emph{band-bottom} reduction:
by Proposition~\ref{prop:wschur},
$\rho_s(c,k)\le\bar\rho(C,h,k,s):=
(s{+}1)W_{s-1}(\mathrm{ext})W_{s+1}(\mathrm{ext})/(sW_s(\mathrm{bal})^2)$
with $\mathrm{ext}=(1^{h-1},C{-}h{+}1)$, an explicit inequality
between two canonical multisets.  Exact maps ($h\le60$, $C\le2000$)
give $\bar\rho\le0.546$ at $k-s=1$ (tracking the derived asymptote
$(2-\tfrac32\lambda)/(2-\lambda)^2\le\tfrac9{16}$,
$\lambda=(s{-}1)/h$) and $\bar\rho\le0.897$ at $k-s=2$, while
$\bar\rho$ crosses $1$ at $k-s=3$ --- the elementary reduction
reaches exactly to where the high regime begins, so the open
domination lemma is confined to $k-s\ge3$.  At rung $1$ the ratio is the quadratic form
$\sum_{i\ne j}w_iw_jK(c_i,c_j)$ in the normalized weights
$w_i=x_i/\sum x$, with the $h$-free kernel
$K(\sigma,\tau)=g\,y(\sigma{+}\tau)/(x(\sigma)x(\tau))$, which is
log-submodular (equivalently, $y$ is log-concave in its argument)
and maximized at extreme pairs that real multisets cannot weight
--- the mechanism of domination is this weight--kernel
anticorrelation, and exact probes rule out every reduction that
decouples a pair from the rest of the multiset.  What the probes
leave standing is a variational program: over the transfer graph
at fixed $(C,h,k)$, all observed tight local maximizers are
either adjacent-two-value or of the
\emph{bulk-plus-giant} form $(a,\dots,a,M)$.  No domination
between the families survives scrutiny --- exact sweeps of the
giant family show it beating the balanced companion deep in the
tight regime (e.g.\ $(1^5,95)$ at $k=46$: $0.778$ against
$0.734$; balanced ratios up to $0.956$ are beaten), while the
ladder itself stays unviolated, with the sweep maximum $0.9865$
attained at the minimal gap-two giant $(6^{99},8)$.  Domination
is also unnecessary: if every tight local maximizer is one of the
two shapes, the rung for all multisets follows from per-family
bounds alone.  Both candidate-maximizer families are certified within the
standing caps: the adjacent-two-value rungs ($7{,}866$ instances)
and the bulk-plus-giant rungs ($1{,}512$ instances, $h\le10$,
$M-a\le8$, $k\le28$, all integer $a$, zero failures).  An exact
hill-climbing census at scale ($h\le24$, $C\le128$, $\sim18$
starts per position) found $76$ distinct local maximizers, every
one in the two families.  The open core is therefore the
first-order \emph{exchange lemma}: every tight local maximizer of
a rung over the transfer graph is adjacent-two-value or
bulk-plus-giant.  Its proof must treat the stationarity system
globally --- the giant's stability is collective (the pair
$(a,M)$ is stable only because the bulk pins the background), the
structural reason the pairwise reductions above fail.  Second,
monotonicity or limit arguments
upgrading the certificate grid to all parameters
(Conjecture~\ref{conj:g2}).  Third, an aggregate analysis of the
symmetric form in the deep band of the uniform case, where the defect
is observed to be negative but termwise and low-order symmetrization
arguments provably fail; the higher ladder rungs, which carry more
slack than the E-core rung in all scans, may help here too.

\section*{AI-use Declaration}

The author used AI tools, including Claude (Anthropic), during proof
exploration, drafting, code development, and audit.  These tools are
not authors.  The human author reviewed the output, takes responsibility
for the manuscript, proofs, verification code and data, and approved the
final version.  All verification scripts and paper-facing logs are
available in the accompanying repository.

\begin{thebibliography}{99}
\bibitem{AMSE87} Y.~Alavi, P.~J. Malde, A.~J. Schwenk, P.~Erd\H{o}s,
  The vertex independence sequence of a graph is not constrained,
  \emph{Congr.\ Numer.} \textbf{58} (1987), 15--23.
\bibitem{ErdosProblems} T.~Bloom, The Erd\H{o}s problems website,
  problem 993, \url{https://www.erdosproblems.com/993}.
\bibitem{BasitGalvin} A.~Basit, D.~Galvin, On the independent set
  sequence of a tree, \emph{Electron. J. Combin.} \textbf{28} (2021),
  no.~3, Paper No.~3.23, \url{https://doi.org/10.37236/9896}.
\bibitem{GalvinHilyard} D.~Galvin, J.~Hilyard, The independent set
  sequence of some families of trees, \emph{Australas. J. Combin.}
  \textbf{70} (2018), no.~2, 236--252.
\bibitem{KadrawiLevit} O.~Kadrawi, V.~E. Levit, The independence
  polynomial of trees is not always log-concave starting from order
  26, \emph{Ars Math. Contemp.} \textbf{25} (2025), no.~4, Paper
  No.~P4.03, \url{https://doi.org/10.26493/1855-3974.3207.2ad}.
\bibitem{GalvinNonLC} D.~Galvin, Trees with non log-concave independent
  set sequences, arXiv:2502.10654.
\bibitem{RamosSun} E.~Ramos, S.~Sun, An AI enhanced approach to the
  tree unimodality conjecture, arXiv:2510.18826.
\bibitem{Li} G.~M.~X. Li, Unimodality of independence polynomials of
  two family of trees, arXiv:2603.03025.
\bibitem{BautistaRamos} C.~Bautista-Ramos, C.~Guill\'en-Galv\'an,
  P.~G\'omez-Salgado, Linear recurrences for non-log-concave
  independence polynomials of trees, arXiv:2603.14204.
\bibitem{Karlin} S.~Karlin, \emph{Total Positivity}, Vol.~I,
  Stanford University Press, 1968.
\bibitem{ShakedShanthikumar} M.~Shaked, J.~G. Shanthikumar,
  \emph{Stochastic Orders}, Springer, 2007.
\bibitem{Liggett} T.~M. Liggett, Ultra logconcave sequences and
  negative dependence, \emph{J.\ Combin.\ Theory Ser.\ A} \textbf{79}
  (1997), 315--325.
\end{thebibliography}

\end{document}
